\chapter{Tutorial: silicon from start to finish}
\label{ch:tutorial-silicon}

This tutorial walks the whole arc of a materials study on one small,
well-understood system: crystalline silicon. You will build the crystal,
relax it, converge a ground state, compute the phonon dispersion, and finish
with the electronic band structure and projected density of states.

Silicon is the right first system precisely because the answers are known.
Every stage below ends with a number you can check, so when something goes
wrong you find out at that stage instead of three stages later.

\begin{calwarn}
GPAW is used throughout. Install it into a Conda environment and point
Calango at it once, in \menu{Edit}\menusep\menu{Preferences}\menusep\menu{Python
\& Environments}. Every wizard then binds to that environment automatically.
Without GPAW you can still follow the tutorial with the \ui{EMT} calculator ---
the workflow is identical --- except for the absence of electronic structure calculations.
\end{calwarn}

\section{Step 1 --- Build the crystal}

Open \menu{Build}\menusep\menu{From Database\dots} and stay on the \ui{Bulk}
tab. It is already set up for this system:

\begin{description}
  \item[\ui{Build from}] \opt{Prototype (ase.build.bulk)}.
  \item[\ui{Formula}] \opt{Si}.
  \item[\ui{Crystal structure}] \opt{diamond}.
  \item[\ui{Lattice constant a}] leave at \opt{0} to accept ASE's tabulated
    value, 5.43\,\AA.
  \item[\ui{Cell shape}] leave unchecked for the two-atom primitive cell.
\end{description}

Click \ui{Build Crystal} and the crystal opens in a new workspace tab. The
\ui{Structure} dock on the left should now read \opt{Si2}, 2 atoms, with
$a = b = c = 3.84$\,\AA{} and $\alpha = \beta = \gamma = 60^\circ$: that is
the rhombohedral primitive cell of the diamond lattice, not an error. The
conventional cubic cell with its 8 atoms and $90^\circ$ angles is what you get
by ticking \ui{Conventional cubic cell}.

\begin{caltip}
Work in the primitive cell. Every stage after this one scales with the number
of atoms --- the phonon step worst of all, since it builds a supercell of
whatever you hand it. Eight atoms instead of two would make that step roughly
an order of magnitude slower for identical physics.
\end{caltip}

\section{Step 2 --- Relax the geometry}

Silicon's tabulated lattice constant is not PBE's lattice constant, so the
first job is to let the calculator find its own minimum.

Open \menu{Simulation}\menusep\menu{Geometry Optimization\dots}. The wizard
runs in four stages; the \ui{Next} button carries you through them.

\begin{enumerate}
  \item \textbf{Calculator \& Execution Environment.} Choose \ui{GPAW}. The
    Conda environment is filled in from Preferences.
  \item \textbf{Calculator Settings.} For this system:
    \begin{itemize}
      \item \ui{Mode} \opt{PW}, \ui{Plane-wave cutoff} \opt{400\,eV}
      \item \ui{XC functional} \opt{PBE}
      \item \ui{k-points} \opt{8 $\times$ 8 $\times$ 8}
    \end{itemize}
  \item \textbf{Optimization Settings.} \ui{Force convergence (fmax)}
    \opt{0.01\,eV/\AA}. Silicon's primitive cell has no free internal
    coordinates by symmetry, so this converges almost immediately.
  \item \textbf{ASE Script Review.} Read the script --- it is the whole job,
    and it is editable. Press \ui{Run}.
\end{enumerate}

The \ui{Processes} dock tracks the run and the \ui{Results} dock plots energy
against step as it goes. When it finishes, the relaxed structure loads back
into the tab and \menu{Results}\menusep\menu{Geometry Optimization Viewer\dots}
shows the energy and force convergence history.

\begin{calnote}
Relaxing \emph{positions} will not change silicon's lattice constant --- BFGS
moves atoms, not the cell. Getting PBE's lattice constant (about 5.47\,\AA,
some 0.7\,\% above experiment, which is PBE behaving normally) means an
energy--volume scan: build several cells with
\menu{Build}\menusep\menu{Supercell\dots} or the structure editor, run a
single point on each, and fit the minimum. This tutorial keeps the tabulated
constant so the later stages stay comparable to published numbers.
\end{calnote}

\section{Step 3 --- Converge the ground state}

\menu{Simulation}\menusep\menu{Single-point Calculation\dots} (\kbd{Ctrl}+\kbd{R})
now does double duty. It reports the total energy, and --- because the
calculator is GPAW --- it writes \file{single\_point.gpw}, the converged
density and wavefunctions.

Use the same calculator settings as Step 2, but raise the k-grid to
\opt{12 $\times$ 12 $\times$ 12}. Run it.

\begin{calnote}
This file is the \emph{baseline} that the later analysis workflows inherit.
Electronic Structure, Optics, MLWF, GW and the charge density difference all
load a completed ground
state and evaluate at fixed density rather than re-converging their own. That
is deliberate: a spectrum computed from a silently different SCF solution than
the one you inspected is not attributable to anything. It also means those
wizards will refuse to open until this step exists --- if one tells you it
needs a baseline, this is the step you skipped.
\end{calnote}

\menu{Results}\menusep\menu{Single-Point Viewer\dots} shows the total energy,
the Fermi level and the maximum residual force.

\section{Step 4 --- Phonon dispersion}

Open \menu{Simulation}\menusep\menu{Phonon\dots}. This wizard has four stages.

\begin{enumerate}
  \item \textbf{Calculator \& Execution Environment.} GPAW again.
  \item \textbf{Calculator Settings.} As before.
  \item \textbf{Phonon Settings.}
    \begin{itemize}
      \item \ui{Displacement $\delta$} \opt{0.01\,\AA} --- small enough to stay
        harmonic, large enough to lift the forces clear of numerical noise.
      \item \ui{Supercell} \opt{3 $\times$ 3 $\times$ 3}. The supercell sets
        how far force constants are sampled, and therefore how much of the
        dispersion is real rather than interpolated. $2\times2\times2$ runs
        faster and already gives the right shape.
      \item \ui{Symmetry-reduced displacements} \textbf{on}. The naive scheme
        costs $6N$ force evaluations; silicon's space group relates almost all
        of them, and spglib (through phonopy) reduces the set by roughly an
        order of magnitude for identical force constants.
      \item \ui{Remove residual forces} \textbf{on}. A relaxation stops at a
        finite \opt{fmax}, so the reference geometry still carries small
        forces; left in, they add a spurious linear term that shows up as
        non-zero acoustic frequencies at $\Gamma$.
      \item \ui{Enforce acoustic sum rule} \textbf{on}.
    \end{itemize}
  \item \textbf{q-Path Definition.} The embedded Brillouin-zone builder
    suggests silicon's conventional path,
    $\Gamma \to X \to W \to K \to \Gamma \to L$. Accept it.
\end{enumerate}

This is the longest step of the tutorial --- it is many force evaluations on a
supercell. When it completes, the \ui{Phonon Viewer} opens on its own.

\subsection{Checking the answer}

Two atoms in the primitive cell give $3 \times 2 = 6$ branches: three acoustic
and three optical. Check these before reading anything else into the plot:

\begin{description}
  \item[Acoustic branches vanish at $\Gamma$.] All three should start at zero.
    Residual frequencies of a few cm$^{-1}$ are normal; tens of cm$^{-1}$ mean
    the acoustic sum rule or the residual-force subtraction is not doing its
    job, or the relaxation was too loose.
  \item[The optical mode at $\Gamma$ sits near 15.5\,THz (520\,cm$^{-1}$).]
    This is silicon's Raman line, one of the best-measured numbers in solid
    state physics. Within a few percent means the calculation is sound.
  \item[No imaginary frequencies.] Plotted as negative, they would mean the
    structure is not at a minimum --- for silicon that indicates a problem with
    the setup, not a real instability.
\end{description}

\menu{Analysis}\menusep\menu{Phonon Thermodynamics\dots} integrates the density
of states into internal energy, free energy, entropy and heat capacity over
0--1000\,K. Two checks come free: $C_V$ must approach the Dulong--Petit limit
$3Nk_B$ at high temperature, and the entropy must go to zero at 0\,K.

\section{Step 5 --- Electronic structure}

Open \menu{Simulation}\menusep\menu{Electronic Structure\dots}. Stage 1 asks
for a \ui{Baseline SCF density} --- the \file{single\_point.gpw} from Step 3.
Select it.

The run is non-self-consistent by construction: it loads that density and
evaluates the bands along the k-path at fixed density. Set the k-path the same
way as for the phonons ($\Gamma \to X \to W \to K \to \Gamma \to L$) and enable
\ui{PDOS} to get the element- and orbital-projected density of states beside
the bands.

When it finishes, the band/PDOS viewer opens: bands as $E - E_F$ against
k-distance on the left, the PDOS sharing the energy axis on the right.

\subsection{Checking the answer}

\begin{description}
  \item[The gap is indirect.] The valence-band maximum sits at $\Gamma$; the
    conduction-band minimum does not --- it lies along $\Gamma \to X$, about
    85\,\% of the way to $X$. A band structure showing silicon's minimum at
    $\Gamma$ is wrong.
  \item[The gap is too small, and that is expected.] PBE gives roughly
    0.6\,eV against an experimental 1.17\,eV. This is the standard
    band-gap underestimation of semilocal DFT, not a convergence failure.
    Raising the cutoff or the k-grid will not fix it, because it is not an
    error in the calculation.
  \item[The PDOS is $sp^3$.] The valence band is a mixture of Si $s$ and $p$
    character, as tetrahedral bonding requires.
\end{description}

\begin{caltip}
To correct the gap rather than accept it, run
\menu{Simulation}\menusep\menu{GW Calculations\dots} on the same baseline.
One-shot $G_0W_0$ replaces the DFT exchange--correlation potential with the
self-energy and opens semiconductor gaps by typically 0.5--2\,eV. The
\ui{GW Viewer} reports the DFT gap, the quasiparticle gap and the
renormalization between them.
\end{caltip}

\section{Where to go next}

The same baseline supports the rest of the analysis suite without recomputing
the ground state:

\begin{itemize}
  \item \menu{Simulation}\menusep\menu{Optics\dots} --- the dielectric function
    $\varepsilon(\omega)$ and the absorption, reflectivity, refractive-index
    and loss spectra derived from it.
  \item The \ui{ELF} checkbox under \ui{Density Exports} in the single-point
    setup --- the covalent bond charge between neighbours, written as
    \file{elf.cube} by the same SCF and rendered as an isosurface in the
    \ui{Volumetric Data} dock. There is no separate ELF module: the field
    comes out of the ground state you already have.
  \item \menu{Analysis}\menusep\menu{Charge Density Difference (CDD)\dots} ---
    $\Delta\rho = \rho_{A+B} - \rho_A - \rho_B$ between two fragments of the
    same run, showing where the charge went when they were put together.
  \item \menu{Simulation}\menusep\menu{Maximally Localized Wannier Functions
    (MLWF)\dots} --- the four $sp^3$ bond orbitals of the diamond lattice.
  \item \menu{Analysis}\menusep\menu{Raman Modes\dots} --- the factor-group
    analysis that labels silicon's $\Gamma$ optical mode $T_{2g}$ and
    Raman-active.
\end{itemize}

Save the workspace with \menu{File}\menusep\menu{Save} before you stop. Job
directories of a saved project are kept alongside the \file{.calproj}, so
every result above stays linked in the \ui{Processes} dock and can be reopened
without recomputation.
